\section{Conclusion}

This paper presented \sys, which explains frame hitches in unmodified Linux
games by partitioning the presenting thread's frame by kernel cause,
resolving its waits online through wake chains carried in task-local
storage, and deciding in the kernel, at frame end, whether the frame is
worth a record. On a ground-truth benchmark it recorded exactly the
over-budget frames, named the injected mechanism, and distinguished a
thread waiting on a stalled worker from one waiting on a busy worker. In ten
case studies its records explained nearly every hitch whose cause was known,
and the failures of an earlier version shaped the design: the waker's
off-CPU history and the record of who held a queued task's CPU each fixed a
case. On rare hitches it was the only one of eight tools to attribute nearly
every hitch and name its culprit, at a few hundred times less output than
full tracing or a flight recorder, for a cost per context switch in the range of existing
scheduler tracers. The case studies also list plainly what the records
cannot yet say: that a sleep was late, which CPU a thread waited for, and
who else was using the disk.
\tbd{[Conclusion on real titles and prevalence.]}

\section*{Acknowledgments}

\todo{Acknowledgments and funding.}

\section*{Availability}

\todo{Repository URL, license, and the scripts that regenerate every table
and figure from raw runs.}
